% Prologos Type System — LaTeX Wrapper for Ott-Generated Rules
%
% Usage: make pdf   (generates prologos-rules.pdf)

\documentclass[11pt]{article}

\usepackage[margin=1in]{geometry}
\usepackage{amsmath,amssymb,amsthm}
\usepackage{mathpartir}
\usepackage{supertabular}
\usepackage{hyperref}
\usepackage{xcolor}

% Ott-generated definitions
\input{prologos-generated.tex}

\title{Prologos Type System\\[0.3em]
\large Inference Rules and Formal Specification}
\author{The Prologos Project - from Zachary Larson}
\date{February 2026}

\begin{document}
\maketitle

\begin{abstract}
This document presents the formal type system specification for Prologos,
a functional-logic language unifying dependent types, session types,
linear types (QTT), logic programming, and propagators. The specification
is generated from the Ott source in \texttt{prologos.ott}.
\end{abstract}

\tableofcontents
\newpage

\section{Syntax}

The syntax of Prologos uses a \emph{locally-nameless} representation:
bound variables are de~Bruijn indices ($\mathit{bvar}~k$), while free
variables are named ($\mathit{fvar}~x$). This avoids $\alpha$-equivalence
issues in the metatheory.

\subsection{Expressions}

\ottgrammartabular{ottexpr\ottinterrule}

\subsection{Multiplicities (QTT)}

The multiplicity semiring $\{0, 1, \omega\}$ governs resource usage:
\begin{itemize}
  \item \textbf{Zero}: erased (compile-time only, 0 runtime uses)
  \item \textbf{One}: linear (exactly 1 runtime use)
  \item \textbf{Omega}: unrestricted (any number of uses)
\end{itemize}

\ottgrammartabular{ottmult\ottinterrule}

\subsection{Universe Levels}

Universes are stratified: $\mathsf{Type}~\ell : \mathsf{Type}~(\mathsf{lsuc}~\ell)$.

\ottgrammartabular{ottlevel\ottinterrule}

\subsection{Session Types}

Session types describe communication protocols on channels:

\ottgrammartabular{ottsession\ottinterrule}

\subsection{Branches}

\ottgrammartabular{ottbranches\ottinterrule}

\subsection{Typing Contexts}

\ottgrammartabular{ottctx\ottinterrule}

\subsection{Usage Contexts (QTT)}

\ottgrammartabular{ottusage\ottinterrule}

\subsection{Processes}

\ottgrammartabular{ottproc\ottinterrule}

\subsection{Channels and Channel Contexts}

\ottgrammartabular{ottchannel\ottinterrule}
\ottgrammartabular{ottchanctx\ottinterrule}

\subsection{Process Branches}

\ottgrammartabular{ottprocbranches\ottinterrule}


\section{Reduction}

Reduction proceeds in two phases: \emph{weak-head normal form} (WHNF) for
type checking, and \emph{full normalization} (NF) for definitional equality.

\subsection{Single-Step WHNF Reduction}

\ottdefnwhnfXXstep{}

\subsection{Weak Head Normal Form (Transitive Closure)}

\ottdefnwhnf{}

\subsection{Definitional Equality}

Two terms are definitionally equal ($e_1 \equiv e_2$) iff their full normal
forms are syntactically identical: $\mathit{nf}(e_1) = \mathit{nf}(e_2)$.


\section{Typing}

Prologos uses a \emph{bidirectional} type system with two modes:
\begin{itemize}
  \item \textbf{Synthesis} ($\Gamma \vdash e \Rightarrow A$): infer the type of $e$
  \item \textbf{Checking} ($\Gamma \vdash e \Leftarrow A$): check $e$ against expected type $A$
\end{itemize}

\subsection{Type Inference (Synthesis)}

\ottdefninfer{}

\subsection{Type Checking}

\ottdefncheck{}


\section{Quantitative Type Theory (QTT)}

QTT extends the type system with usage tracking. Each judgment returns
a \emph{usage context} $\mathcal{U}$ recording how many times each
variable was used:

\begin{itemize}
  \item \textbf{QTT Inference}: $\Gamma \vdash e \Rightarrow_Q A \;;\; \mathcal{U}$
  \item \textbf{QTT Checking}: $\Gamma \vdash e \Leftarrow_Q A \;;\; \mathcal{U}$
\end{itemize}

Usage contexts compose via pointwise semiring operations:
$\mathcal{U}_1 + \mathcal{U}_2$ (addition) and
$q \cdot \mathcal{U}$ (scalar multiplication).

\subsection{QTT Inference}

\ottdefninferQ{}

\subsection{QTT Checking}

\ottdefncheckQ{}


\section{Session Types}

Session types describe bidirectional communication protocols. The key
operation is \emph{duality}: if one endpoint has session $S$, the other
has $\overline{S}$.

\subsection{Duality}

\ottdefndual{}


\section{Process Typing}

Processes are typed against channel contexts $\Delta$ mapping channels
to their session types:

$$\Gamma ; \Delta \vdash P$$

\ottdefntypeProc{}


\section{Metatheoretic Properties}

The following properties are verified by the PLT Redex formalization
(\texttt{racket/prologos/redex/properties.rkt}) via randomized testing:

\begin{enumerate}
  \item \textbf{Type Preservation}: If $\vdash e \Rightarrow T$ and $e \to e'$,
        then $\vdash e' \Rightarrow T'$ with $T \equiv T'$.
  \item \textbf{Progress}: If $\vdash e \Rightarrow T$ (closed), then $e$ is a
        value or $e \to e'$.
  \item \textbf{Determinism}: WHNF reduction is deterministic.
  \item \textbf{Conversion is an Equivalence}: Reflexive, symmetric, transitive.
  \item \textbf{Shift Identity}: $\mathit{shift}(0, c, e) = e$.
  \item \textbf{Duality Involution}: $\overline{\overline{S}} = S$.
  \item \textbf{QTT Soundness}: QTT checking refines standard type checking.
\end{enumerate}

\end{document}
